% ============================================================
% Related Work — CHI 2027 投稿草稿（中文原文，待翻译成英文）
% 生成自: relatedwork/Related work 成稿.md
% 所有 \cite key 与 references.bib 一一对应（2026-08-02 已做引用核验）
% ============================================================
\section{Related Work}

本研究的核心任务是：从海量众包灾害新闻中自动筛选出可被 VR 重建的具身场景叙事。这个任务横跨了三条 HCI 研究脉络：一是沉浸式新闻与新闻叙事 VR，追问"如何让受众以第一人称经历新闻事件"；二是 VR 具身感与共情设计，追问"身体在场如何转化为情感共鸣"；三是弱势群体的 VR 可达性与表征，追问"如何让 VR 体验对所有人都可用、可感、可被呈现"。下面先按照这三条脉络梳理已有工作的共识与局限，再整合指出研究空白与本文贡献。

\subsection{Immersive Journalism and Narrative VR}

\subsubsection{沉浸式新闻的定义与理论}

沉浸式新闻（Immersive Journalism）的概念由 de la Peña 等人于 2010 年正式提出，定义为"以让人们对新闻故事中的事件或情境获得第一人称体验的形式来生产新闻"\cite{delapena2010immersive}。该定义界定了沉浸式新闻的三个关键特征：非虚构性——内容基于真实新闻事件；第一人称在场——用户以自我为中心的视角存在于虚拟空间中；时间同步性——用户以实时方式经历事件进程而非事后回看。这一定义奠定了沉浸式新闻作为一种新闻呈现范式的基础：受众从信息接收者转变为事件的空间参与者。

然而，仅凭借这一定义和理论描述，尚不足以支撑从大规模灾害新闻文本到 VR 叙事体验的完整转换。现有研究缺少一套从文本到 VR 场景的可落地、可执行的生成方案，同时也缺少在生成过程中对叙事与交互机制的系统化设计方法。如何将新闻正文中包含的具身细节（空间描述、身体动作、感知体验）有效地转化为 VR 中可交互的叙事场景，仍是尚未被充分解决的问题。

\subsubsection{效果检验：从"默认共情"到"可设计的共情"}

在效果检验方向，Feeling the News? 的对照实验提出了一个与"第一人称新闻"理想形成对照的发现。实验数据显示，被动观看 360° 视频在情感参与度上显著低于主动交互式 VR 体验；叙事结构（线性 vs. 分支）对受众的认知共情也有显著影响，分支叙事组在事后测试中对新闻事件细节的回忆更准确。作者的结论是："Immersive journalism (IJ) is often assumed to be inherently emotion-inducing... the audience responses depend on interaction and narrative design"\cite{feelingnews2023}。上述结果表明，沉浸本身只是内容传输的媒介条件，受众的情感反应取决于可被调控的叙事编排与交互机制设计。

Herrera 和 Bailenson 的大规模实验（超过 400 名参与者）从另一个角度支持了这一观点。他们比较了传统视角转换任务和 VR 视角转换对长期共情的影响，发现 VR 组在观点采择上的提升显著高于传统组，且效应在一周后仍可测量\cite{herrera2018building}。这说明沉浸式环境下的共情效应并非自然发生，而是需要经过设计才能产生可测量的持久影响。

\subsubsection{设计框架与产业瓶颈}

上述发现自然引出一个问题：\textbf{如何构造有效的叙事和交互机制以激发受众的共情反应？} 近年来，多个研究从不同层面回应了这一问题。在评估层面，Beyond "Being There"（CHI 2026）将 VR 叙事参与度分解为认知（空间记忆、叙事理解）、情感（情绪效价、共情反应）、行为（交互深度、路径选择）三个维度，并提出了对应的测量方法\cite{beyondbeingthere2026}。在范式层面，Storyliving 提出应将 VR 理解为"叙事生活"（storyliving）——用户通过身体行为主动建构意义，而非被动接收作者预设的叙事\cite{storyliving2022}。在工具层面，RécitKit 提供了一套空间叙事设计套件，包含时空场景编辑器、注意力引导系统和评估模块，可将文本描述转化为 VR 中的空间叙事单元\cite{recitkit2025}。

从产业实践的角度观察，Embracing Immersive Journalism 的调研揭示了沉浸式新闻面临的产业化瓶颈：典型制作预算在 5 万到 50 万美元之间，此外还存在人才缺口、分发困难和 ROI 不明确等问题\cite{embracing2024}。现有的生产模式（现场拍摄 + 人工制作）成本过高，难以支持大规模的新闻叙事内容生产。这也是现有沉浸式新闻工作主要依赖少量被人工精选的素材（如知名通讯社稿件）的原因之一。

\subsection{Embodiment and Empathy in VR}

VR 共情研究的理论根基在于具身感——身体在虚拟空间中的存在感如何影响情感反应。Kilteni 等人给出了具身感（Sense of Embodiment, SoE）的操作性定义，将其分解为三个子要素：自我定位（self-location），即感知到自己位于虚拟身体内部的体验；身体所有权（body ownership），即虚拟身体属于自己的感觉；能动性（agency），即对虚拟身体动作有控制感的体验\cite{kilteni2012embodiment}。该工作（Google Scholar 引用超过 2,400 次）已成为后续 VR 具身研究的基础框架。研究者通过橡胶手错觉的 VR 版实验发现，视觉-触觉一致性可以显著增强对身体所有权的感知，且此效应可推广至整个虚拟身体；通过延迟运动反馈实验（100ms/300ms/500ms 延迟量），发现 100ms 以内的延迟对能动感影响较小，300ms 以上则显著破坏能动性。

SoE 三要素构成了本文筛选标准中"具身场景可重建性"的理论根系。空间细节对应自我定位（新闻报道提供了多少可被 VR 重建的楼层、走廊、楼梯信息？），身体动作对应身体所有权（报道是否描写了当事人在火灾中的具体身体行为——扶着墙走、匍匐前进？），决策描写对应能动性（报道是否呈现了当事人的选择困境——不乘电梯、改走消防通道？）。我们在筛选新闻时，正是用这三条标准来判定一篇正文是否具备 VR 重建价值。

早于 Kilteni 的工作，Slater 和 Wilbur 在 1997 年就系统定义了沉浸式虚拟环境的临场感（presence）框架，将"在场"区分为技术层面的沉浸（immersion）和心理层面的临场感（presence）\cite{slater1997five}。这一区分至今仍有理论价值：沉浸是 VR 系统的技术参数（视场角、追踪精度、显示分辨率），临场感是用户的心理状态（"我好像真的在那里"），而共情是更高层次的情感结果。这三者不能混为一谈——Feeling the News? 的发现（被动观看即使在场感高，情感提升也不显著）\cite{feelingnews2023}正是在实践中验证了这一理论区分。

在设计路径上，Jiang 等人提出了一个 VR 共情设计框架（EVRD），发表于 IJHCS 2025。该工作采用实践导向的研究方法（practice-led design），提出了一个反直觉的设计策略——"doomed-to-fail"（注定失败）。传统 VR 叙事往往设计用户"成功救援"的情节，但 EVRD 的实验结果显示，让用户在 VR 中体验"无法成功"的情境（如无论怎么努力都无法打开那扇门），对认知共情的促进作用反而更显著——用户通过亲身经历理解了障碍者所处困境的结构性原因，而非停留在情感同情层面\cite{evrd2025}。

Ahn 和 Bailenson 的经典实验从另一个角度验证了具身体验的力量：他们让参与者在一款沉浸式虚拟环境中"砍伐"一棵虚拟树，或者仅仅观看别人砍树的视频。结果发现，具有具身体验的参与者在事后的环保行为和态度改变上显著强于仅观看视频的对照组\cite{ahn2013embodied}\cite{ahn2014short}。这说明具身体验不仅影响情感，还能驱动实际的行为改变——这在论文方法学上意味着：我们对 VR 共情效果的评估不能只测情绪反应，还应该测量行为层面的变化。

在更宏观的层面，近年来也出现了几篇整合性的综述。Cultivating Empathy Through Narratives in VR（2024）系统回顾了叙事驱动的 VR 共情研究，将现有研究按照"叙事类型×共情维度"进行了分类\cite{cultivating2024}。一篇发表在 Virtual Reality 期刊上的范围综述则更为全面地梳理了 VR 共情干预研究的全景，识别了当前研究在方法论一致性上的不足\cite{vrscoping2024}。这些综述为本文的贡献定位提供了参照：相比于"从零设计一个 VR 共情体验"（大多数现有工作），我们关注的是一个前置但被忽视的问题——"从海量文本中筛选出哪些素材具备 VR 共情重建潜力"。

\subsection{VR for Vulnerable Groups: Accessibility, Representation, and Ethics}

将 VR 用于弱势群体体验与共情培育，近年在 CHI 会议上涌现了一系列值得关注的工作。这些工作可以从"怎么做"（方法论）、"怎么理解"（理论框架）、"怎么设计"（设计指南）和"怎么批评"（伦理约束）四个层面来理解。

在方法论层面，Shocking Realities（DIS 2025）通过 VR 恐怖游戏将轮椅用户的日常挑战嵌入恐怖游戏机制之中，设计了四种具体的场景关卡：狭窄通道、坡道角度（超出 ADA 标准的斜坡）、到达高度（安装在站姿者可达高度而非轮椅使用者可达高度的消防报警按钮）、以及地板障碍（地面上散布碎片和障碍物，轮椅只能绕过而非跨过）。用户测试显示，体验过的参与者对轮椅可达性问题的认知显著提升，而且这种认知提升不是通过"怜悯"而是通过"体验到了障碍确实存在"来实现的\cite{shockingrealities2025}。这个方法论突破在于它避免了传统说教式共情的陷阱——不是告诉用户"轮椅用户很困难"，而是让用户"在游戏中自己发现困难"。我们采用同一条路径：不让用户被告知"听障者在火灾中很危险"，而是让用户以听障者的化身经历只有声音警报而没有闪光信号的环境。

另一组对照实验来自视觉障碍模拟的研究。该实验让 35 名大学生以第一人称体验四种视觉受限场景（近视模糊、中心视野缺失、周边视野缺失、完全黑暗），并在事后使用 IRI 量表测量共情变化。结果显示，实验组在观点采择和共情关注维度上有显著的统计学提升（p < 0.01, d = 0.45），且一周后的追踪测试中效应仍部分保留\cite{visualimpairment2025}。The Machine to Be Another 则从另一个角度切入具身共情：通过实时摄像机画面交换让两个参与者"拥有"对方的身体，后续 fMRI 研究显示体验者的前额叶皮层激活模式在内群体-外群体区分任务中发生改变\cite{machinetobeanother}。

在理论框架层面，An Equitable Experience?（ACM TACCESS）通过系统性文献综述，识别出了 HCI 文献中关于 VR 残障可达性的三种概念化方式：技术可达性（硬件能否使用，如控制器是否可单手操作）、体验可达性（体验质量是否与非残障用户相当，如盲人用户能否通过空间音频获得类似的场景认知）、表征可达性（残障身份是否在 VR 内容中得到准确尊重多样化的呈现，如化身设计是否避免刻板印象）。该研究还指出，当前大多数工作集中在技术可达性层面，表征可达性最少被触及\cite{equitable2025}。这直接提醒本文：我们的生成式城市 VR 必须在三个层面同时确保可达性。

在设计指南层面，Inclusive Avatar Guidelines（CHI 2025）提出了四条残障化身设计原则：提供选项而非强制（让用户自己决定是否在化身中展示辅助器具）、避免唯一表征（同一残障类别应有不同年龄、性别、风格的化身变体）、功能对应而非视觉装饰（辅助器具在 VR 中应与真实交互功能对应）、与真实社区协同设计\cite{inclusiveavatar2025}。这些指南为本文的"城市人格"化身设计提供了直接参考——用户在 VR 城市中可能以轮椅使用者或听障者的身份体验火灾逃生，其残障表征需要遵循上述指南以避免刻板印象。

LegacySphere（CHI 2024）提供了一个整合了具身、角色扮演和叙事三要素的 VR 视角转换系统，用于跨代际沟通。其核心发现是：具身体验（embodiment）与叙事体验（narrative）之间存在协同效应——两者结合比单独使用任何一个都能更有效地促进共情\cite{legacysphere2024}。这个发现对本文有直接的方法论意义：我们的 VR 城市叙事不止是"让用户看一个火灾新闻"（叙事），也不只是"让用户以第一人称视角在场景中移动"（具身），而是将两者结合——用户以某弱势群体角色的具身视角，在生成式城市中经历来自真实新闻的叙事事件。

在伦理约束层面，Tigwell（CHI 2021）通过调研 92 名设计师和 17 名盲人/低视力用户，对失能模拟的伦理问题提出了系统批评：如果只呈现"缺陷"而不揭示"环境限制"，模拟可能激起怜悯或恐惧而非真正的共情\cite{tigwell2021}。这项 CHI 高引工作已经在方法论层面被上述 DIS 2025 的工作所回应——Shocking Realities 通过"环境失败"机制而非"缺陷展示"来设计体验\cite{shockingrealities2025}，Equitable Experience 则从理论层面将"表征可达性"提升到了与技术可达性同等重要的地位\cite{equitable2025}。本文的设计路径也遵循这一伦理约束：我们不孤立地模拟"听障"或"行动不便"，而是在具体的火灾情境中呈现"身体 × 环境耦合"的困境——例如，在只有声音警报的建筑中，听障者无法收到逃生信号——让用户理解到，问题不在于"听障者的缺陷"，而在于"环境没有为所有人设计"。

\subsection{Research Gap and Contribution}

综合以上三条脉络，可以看到一个共同的研究空白贯穿始终。

现有工作有一个未经言明的隐含前提：\textbf{好内容靠人工精制}。无论是沉浸式新闻的内容来源（依赖知名通讯社的已发表稿件\cite{delapena2010immersive}和昂贵的现场制作\cite{embracing2024}），还是 VR 共情模拟的内容设计（研究人员精心设计实验脚本\cite{shockingrealities2025}\cite{visualimpairment2025}\cite{legacysphere2024}），它们默认的内容供给路径都是"手工"的——人找内容、人挑内容、人做内容。HCI 文献关注的焦点是"如何评估已成型的 VR 体验"\cite{beyondbeingthere2026}和"如何设计可用的 VR 系统"\cite{recitkit2025}，但\textbf{从大规模噪声语料中自动筛选出 VR 可重建的场景}这一前置问题，几乎无人系统化处理。PRISMA 提供了系统综述的报告标准，但 PRISMA 面向的是学术文献检索，而非面向 VR 场景重建的新闻正文筛选\cite{page2021prisma}。

更进一步，已有的弱势群体 VR 共情体验都是\textbf{单次、封闭、预制}的——用户"玩完"一个关卡\cite{shockingrealities2025}或"完成"一次模拟\cite{visualimpairment2025}后体验即结束。缺乏将弱势群体视角嵌入一个\textbf{持续生成的城市空间}的框架，使用户可以在不同条件下反复遭遇同一个弱势群体角色的身体限制。

本文尝试同时回应上述两个不足。在内容来源上，提出了一种人机协作的语料筛选方法——以 103 条人类金标准驱动对 8,180 条新闻正文的筛选，误留率 0%、误删率 0.97%、自动剔除率 87.6%，从中提取出 151 条具备 VR 重建条件的正文。该方法本身构成一种可迁移的 HCI 方法学工具，尝试填补"从文本到 VR 场景"的筛选方法论空缺。在体验形式上，将筛选出的正文嵌入一个持续生成的城市叙事中，使用户可以反复以不同弱势群体的化身经历身体 × 环境耦合的共情叙事。该框架的理论基础来自 SoE 三要素\cite{kilteni2012embodiment}，评估方法参考 Beyond "Being There" 的三维参与度框架\cite{beyondbeingthere2026}，叙事范式借鉴 Storyliving 的"叙事生活"理念\cite{storyliving2022}，弱势群体相关的方法论和伦理约束则来自 Shocking Realities\cite{shockingrealities2025}、Equitable Experience\cite{equitable2025}和 Tigwell 的研究\cite{tigwell2021}。
